\documentclass[11pt]{article}
\usepackage[margin=1in]{geometry}
\usepackage{amsmath,amssymb,amsthm,booktabs}
\usepackage[hidelinks]{hyperref}
\usepackage{xurl}
\setlength{\emergencystretch}{2em}
\newtheorem{lemma}{Lemma}
\newtheorem{theorem}[lemma]{Theorem}
\newcommand{\F}{\mathbb F}
\newcommand{\mr}{\operatorname{mr}}
\newcommand{\rk}{\operatorname{rank}}
\title{Three sharp meridional bounds from a binary affine quotient}
\author{Written proof and exact finite certificates}
\date{17 November 2026}
\begin{document}\maketitle

The knots $13n_{4641},13n_{4642},13n_{4645}$ have meridional rank
and bridge number four. Explicit binary matrices give epimorphisms
from their knot groups to one finite affine group of order $161280$.
That group has ordinary generating rank two, while its meridian
conjugacy class requires four generators. Literal matrices, projected
Schreier kernels, planar-map certificates and four-seed traces support
the proof. A branched-cover homology calculation is a separate optional
control, rather than a premise of the direct finite-group bound.

\section{The affine group and its distinguished class}

Let $A_7$ act on the seven points $0,\ldots,6$, and put
\[
 W=\{(w_0,\ldots,w_6)\in\F_2^7:\textstyle\sum_iw_i=0\},
 \qquad Q=W\rtimes A_7.
\]
Write $P_pe_i=e_{p(i)}$, with multiplication
\[
                  (p,v)(q,w)=(pq,v+P_pw).
\]
The homogeneous action on $V=W\oplus\F_2$ has dimension seven:
\[
                  (p,v)(w,s)=(P_pw+sv,s).
\]
In the basis $e_j+e_6$ for $W$, $0\le j<6$, its matrix is
\begin{equation}\label{eq:matrix}
 L(p,v)=\begin{pmatrix}
 (\mathbf1_{p(j)=i}+\mathbf1_{p(6)=i})_{0\le i,j<6}
       &(v_i)_{0\le i<6}\\0&1
 \end{pmatrix}.
\end{equation}
Every entry is taken modulo two. Let $C$ be the conjugacy class in
$Q$ of $(\sigma,0)$, where $\sigma$ is a double transposition in $A_7$.

\begin{lemma}\label{lem:fixed}
If a group $H$ acts on a finite-dimensional vector space $E$ and is
generated by $k$ elements of displacement rank at most $\delta$,
then
\[
                     \dim E-\dim E^H\le k\delta.
\]
\end{lemma}
\begin{proof}
The common fixed space is the intersection of the kernels of the
$k$ displacement maps. The codimension of an intersection is at
most the sum of the individual codimensions.
\end{proof}

\begin{theorem}\label{thm:group}
The group $Q$ has order $2^6\cdot2520=161280$ and ordinary
generating rank two. Its class $C$ has $420$ elements and generating
rank exactly four. Every element of $C$ has displacement rank two
on $V$, and $V^Q=0$.
\end{theorem}
\begin{proof}
The even permutation group has order $7!/2=2520$, and $W$ has
$2^6$ elements. A double transposition has three fixed points and
two moving pairs. The image of $I-P_\sigma$ on $\F_2^7$ is the
two-dimensional span of the pair differences, which here are pair
sums. It lies in $W$. Its image on $W$ is the same: any conjugating
potential can have its coordinate sum adjusted at a fixed sheet
without changing its displacement. Thus
\[
 \operatorname{im}(I-P_\sigma)|_W
          =\operatorname{im}(I-P_\sigma)|_{\F_2^7},\qquad\rk=2.
\]
There are $\binom74\cdot3=105$ double transpositions.
The finite conjugation check verifies that they form one $A_7$ class.
Above each there are exactly four translation conjugates of the
split element, so $|C|=105\cdot4=420$. This is the split conjugacy
class; arbitrary involutory affine lifts are not being identified with it.
By conjugacy and \eqref{eq:matrix}, every element of $C$ has
displacement rank two.

All translations belong to $Q$. A vector $(w,s)$ fixed by all
translations has $s=0$. Transitivity leaves only constant $w$.
A constant vector in $W$ is zero, since the degree seven is odd.
Hence $V^Q=0$, and Lemma~\ref{lem:fixed} gives $7\le2k$ for any
$C$-generating family. Its size is therefore at least four.

For an explicit attaining quartet, set
\[
\begin{array}{ll}
 x_0=((3\,4)(5\,6),e_3+e_4),&
 x_1=((3\,5)(4\,6),e_4+e_6),\\
 x_2=((2\,3)(4\,6),0),&
 x_3=((0\,3)(1\,4),0).
\end{array}
\]
Each belongs to $C$. The exact projected enumeration gives $2520$
elements. For each projected element $p$, choose a lifted word $t_p$
in these generators, with $t_1=1$. For every generator $s$, the loop
\[
                         t_p s t_{ps}^{-1}
\]
is a pure translation. These loops generate the full projection
kernel: inserting the $t_p$ along a kernel word telescopes it into
such loops and their inverses. A negative generator edge is the
inverse of the corresponding positive edge at its other endpoint.
Thus it suffices to use all $2520\cdot4$ positive edges.
The kernel is an elementary abelian binary group, so its exact size
is $2^r$, where $r$ is the row rank of these translation vectors.
The checker obtains $r=6$. Consequently the quartet contains all of
$W$ and projects onto $A_7$, so it generates $Q$.

Finally take $a=x_0x_2$ and $b=x_1x_3$.
Their projected enumeration again has $2520$ elements, and their
$2520\cdot2$ Schreier translations again have rank six. Thus
$\langle a,b\rangle=Q$, without enumerating its $161280$ affine
elements. The nonabelian quotient $A_7$ excludes ordinary rank one.
The ordinary rank is exactly two, and the quartet proves the asserted
upper bound four for $C$-generation.
\end{proof}

\section{Three named knot-group epimorphisms}

\begin{theorem}\label{thm:sharp}
For each
\[
                    K\in\{13n_{4641},13n_{4642},13n_{4645}\},
\]
there is an epimorphism $\pi_1(S^3\setminus K)\twoheadrightarrow Q$
taking meridians into $C$. Each knot has $\mr(K)=b(K)=4$.
\end{theorem}
\begin{proof}
The retained input gives a literal thirteen-crossing PD and thirteen
pairs $(p_i,v_i)$ for each name. Join the edge labels through PD
overpasses in slots one and three, numbering slots from zero.
Sort the resulting connected edge-label components lexicographically
and number them $0,\ldots,12$. These are the Wirtinger arcs.
For a crossing $(o,a,b)$, the first arc is the over-arc and the
other two are under-arcs.

The checker reconstructs each $L_i=L(p_i,v_i)$ from
\eqref{eq:matrix} and compares it with its separately retained literal
matrix. It checks every identity
\[
               L_i^2=I_7,\qquad\rk(L_i-I_7)=2,\qquad
               L_oL_aL_o=L_b.
\]
The last equation satisfies either signed Wirtinger convention,
because all matrices are involutions. They therefore define a
homomorphism from the literal knot group.
Each $p_i$ is a double transposition and $v_i\in W$.
The rank-two condition forces its translation column into
$\operatorname{im}(I-P_{p_i})|_W$, since the permutation displacement
already has rank two. Thus every meridian image belongs to $C$.

The four seed permutations generate $A_7$; their projected Schreier
translations span the full six-dimensional $W$. The images therefore
generate $Q$. For $13n_{4641}$ they are the displayed quartet of
Theorem~\ref{thm:group}. In all three cases the two seed words
$x_0x_2,x_1x_3$ also generate $Q$, as checked by the same exact
projection/kernel calculation.
The common fixed space is zero and every meridian has displacement
rank two. Lemma~\ref{lem:fixed}, or the class-generating rank in
Theorem~\ref{thm:group}, gives $\mr(K)\ge4$.

Each retained four-seed trace reaches every arc using only a known
over-arc and a known under-arc. The published equality between bridge
number and minimum Wirtinger number \cite[Definition~1.2 and
Theorem~1.3]{bkvv} gives $b(K)\le4$. The elementary bridge
presentation gives $\mr(K)\le b(K)$, proving equality throughout.

The literal named-diagram bindings are exact crossing-preserving
dart bijections to retained independent census PDs. They check edge
gluing, vertex rotation and over/under data, allowing a global mirror
or orientation reversal. These operations preserve both ranks and
the involutory meridian condition. The named census is an explicit
trusted import; its name is not inferred from the matrix calculation.
\end{proof}

\section{Reproducibility and the homological control}

Run the standard-library checker from any directory:
\begin{verbatim}
python3 verify_a7_binary_certificates.py
\end{verbatim}
It checks the saved PDs, sphere embeddings, one-component condition,
census dart maps, all matrices and crossings, four-seed propagation,
projected Schreier kernels, ordinary pairs and the $105$-element
projected involution class. It rejects corrupted identity matrices,
seed traces, census maps and ordinary-pair words. The necessary input
is \path{a7_binary_inputs.json}; the result is
\path{a7_binary_verification.json}.

The optional \texttt{--census} mode requires the pinned
\texttt{snappy\_15\_knots}~1.2.1 and Spherogram~2.4.1 environment.
It rereads each named PD and DT and verifies the exact database hash
and byte length. The default checks only the retained named records
and their map certificates.

The optional \texttt{--cover-homology} control reconstructs the full
seven-sheeted branched-cover complex over $\F_2$. It retains all
$91$ directed generator edges, contracts six tree edges, and includes
all $91$ lifted crossing cells and all $65$ meridian-cycle fillings,
including fixed sheets. The resulting $156\times85$ relation matrix
has rank $84$ in every case, so
\[
                       \dim_{\F_2}H_1(Y;\F_2)=1.
\]
After the cycle fillings, a reverse lifted generator edge equals its
forward replacement; their difference is a filling row. Hence using
involutory crossing words gives the same filled row span for either
signed Wirtinger convention. This homology agrees with the cover
bound $\lceil(7-1+1)/2\rceil=4$ and supplies an independent control
of the direct affine result. It is not an indispensable premise.

The general fixed-space argument and conversion of parabolic
cocycles into finite affine representations are not asserted new.
The finite images, their exact orders and the difference between
ordinary and meridian-class generation are additional checked
information. Classification of the double branched covers remains
unresolved here. No worldwide priority, exclusion from every known
MRC family, formal verification or external specialist review is
claimed.

\begin{thebibliography}{9}
\bibitem{bkvv}
R. Blair, A. Kjuchukova, R. Velazquez and P. Villanueva,
Wirtinger systems of generators of knot groups,
Commun. Anal. Geom. 28 (2020), 243--262,
Definition~1.2 and Theorem~1.3, printed p.~246.
The published primary theorem and exact coloring move are imported;
the full theorem is not independently reproved here.
\url{https://archive.intlpress.com/site/pub/files/_fulltext/journals/cag/2020/0028/0002/CAG-2020-0028-0002-a002.pdf}.
\end{thebibliography}
\end{document}
